\chapter{Utility and Indifference Curves}
\label{chap:utility}

\section{Learning objectives}

After completing this chapter, you should be able to:
\begin{itemize}
  \item explain how utility represents preferences;
  \item distinguish ordinal from cardinal utility;
  \item draw and interpret indifference curves;
  \item calculate marginal utility and the marginal rate of substitution; and
  \item recognise the shapes of indifference curves for common preferences.
\end{itemize}

\section{Utility as a representation of preferences}

Preferences rank consumption bundles. A \keyterm{utility function} assigns a number to each bundle in a way that represents this ranking. If $u(A)>u(B)$, then the consumer strictly prefers $A$ to $B$; if $u(A)=u(B)$, the consumer is indifferent between them.

For example, the utility function
\[
u(x,y)=xy
\]
represents a consumer who values both goods. The bundle $(2,6)$ gives utility $12$, while $(3,5)$ gives utility $15$, so the latter is preferred. The numerical values do not measure happiness in an objective sense. Their purpose is to preserve the ordering of bundles.

\section{Ordinal and cardinal utility}

In standard consumer theory, utility is \keyterm{ordinal}. Only the ranking matters, not the size of utility differences. If $u(x,y)=xy$ represents a consumer's preferences, then $v(x,y)=100xy+7$ represents exactly the same preferences: both functions rank all bundles in the same order.

This is different from \keyterm{cardinal utility}, in which numerical differences have an independent meaning. In an ordinal framework, we cannot say that a bundle giving utility $20$ makes a consumer twice as well off as a bundle giving utility $10$. We can say only that the former is preferred.

The ordinal interpretation is important in welfare analysis. Comparing utility numbers across different individuals requires additional ethical assumptions. Economics can describe choices and trade-offs, but interpersonal comparisons of well-being cannot be read mechanically from a utility function.

\section{Indifference curves}

An \keyterm{indifference curve} is the set of all bundles that give a consumer the same utility. For $u(x,y)=xy$, the indifference curve associated with utility level $\bar{u}$ is
\[
xy=\bar{u},
\]
or equivalently
\[
y=\frac{\bar{u}}{x}.
\]

For instance, the bundles $(2,6)$, $(3,4)$, and $(4,3)$ all lie on the indifference curve $xy=12$. The consumer is indifferent among them because each provides the same utility.

Under the usual assumptions of monotonicity and convexity, indifference curves have three important properties:
\begin{enumerate}
  \item they slope downward: receiving less of one good requires more of the other to maintain utility;
  \item higher curves represent higher utility; and
  \item they do not cross: crossing curves would imply contradictory rankings.
\end{enumerate}

\section{Marginal utility}

The \keyterm{marginal utility} of a good is the additional utility associated with a small increase in its quantity, holding the quantity of the other good fixed. For a two-good utility function $u(x,y)$,
\[
MU_x=\frac{\partial u}{\partial x}
\qquad\text{and}\qquad
MU_y=\frac{\partial u}{\partial y}.
\]

For $u(x,y)=xy$,
\[
MU_x=y
\qquad\text{and}\qquad
MU_y=x.
\]
An additional unit of $x$ is more valuable when the consumer has more $y$, because the two goods complement each other in the utility function.

\section{The marginal rate of substitution}

The \keyterm{marginal rate of substitution} of $x$ for $y$, written $MRS_{x,y}$, measures how much $y$ a consumer is willing to give up for one additional unit of $x$ while remaining on the same indifference curve. It is the absolute value of the slope of an indifference curve:
\[
MRS_{x,y}=\frac{MU_x}{MU_y}.
\]

For $u(x,y)=xy$,
\[
MRS_{x,y}=\frac{y}{x}.
\]
At a bundle with much $y$ and little $x$, the consumer is willing to give up a relatively large amount of $y$ for another unit of $x$. As $x$ becomes more abundant and $y$ scarcer, the willingness to substitute falls. This is the \keyterm{diminishing marginal rate of substitution}, which produces bowed-in indifference curves.

\section{Common preference representations}

\subsection{Cobb--Douglas preferences}

The utility function
\[
u(x,y)=x^\alpha y^\beta,
\]
where $\alpha>0$ and $\beta>0$, is called \keyterm{Cobb--Douglas}. It represents smooth, downward-sloping, convex indifference curves. Its marginal rate of substitution is
\[
MRS_{x,y}=\frac{\alpha}{\beta}\frac{y}{x}.
\]
The exponents influence the relative importance of the two goods.

\subsection{Perfect substitutes}

If two goods are perfect substitutes, the consumer is willing to replace one with the other at a constant rate. A utility function is
\[
u(x,y)=x+2y.
\]
Its indifference curves are straight lines. Here, one additional unit of $y$ gives the same utility as two additional units of $x$.

\subsection{Perfect complements}

If goods are valuable only when consumed in fixed proportions, they are perfect complements. For example,
\[
u(x,y)=\min\{x,y\}
\]
can represent left and right shoes. The indifference curves are L-shaped. Additional left shoes without additional right shoes do not increase utility.

\subsection{Quasilinear preferences}

The utility function
\[
u(x,y)=v(x)+y
\]
is called \keyterm{quasilinear}. Its marginal utility of $y$ is constant, while the marginal utility of $x$ may diminish. Such preferences are useful in analysing choices where one good is a general ``outside good,'' such as money spent on other consumption.

\section{What indifference curves do not say}

Indifference curves are not psychological maps. They do not imply that consumers consciously calculate marginal utilities, and they do not show how preferences were formed. They provide a compact way to represent observed or hypothetical rankings. Their value lies in showing the trade-offs implied by those rankings.

\section*{Chapter Review}

\subsection*{Key concepts}
\begin{itemize}
  \item \textbf{Utility function}: a numerical representation of a preference ranking.
  \item \textbf{Ordinal utility}: utility numbers rank bundles but do not measure well-being in absolute units.
  \item \textbf{Indifference curve}: the set of bundles giving the same utility.
  \item \textbf{Marginal utility}: the change in utility from a small increase in one good.
  \item \textbf{Marginal rate of substitution}: the trade-off a consumer accepts while keeping utility unchanged.
  \item \textbf{Diminishing MRS}: willingness to substitute falls as a consumer has more of one good relative to the other.
\end{itemize}

\subsection*{Key equations}
\[
MU_x=\frac{\partial u}{\partial x}, \qquad
MU_y=\frac{\partial u}{\partial y}, \qquad
MRS_{x,y}=\frac{MU_x}{MU_y}.
\]

\section*{Exercises}

\subsection*{Conceptual questions}
\begin{enumerate}[label=\arabic*.]
  \item What does it mean for utility to be ordinal? Why can we replace $u(x,y)$ with $100u(x,y)+7$ without changing preferences?
  \item State and explain the three standard properties of indifference curves.
  \item Define marginal utility and the marginal rate of substitution.
  \item Explain why indifference curves do not normally cross.
  \item Give one example each of perfect substitutes and perfect complements.
\end{enumerate}

\subsection*{Graphical exercises}
\begin{enumerate}[resume,label=\arabic*.]
  \item Draw the indifference curves $u=6$, $u=12$, and $u=18$ for $u(x,y)=xy$. Mark the bundles $(2,3)$, $(3,4)$, and $(6,3)$.
  \item Draw representative indifference curves for perfect substitutes and perfect complements. Explain their different shapes.
  \item For $u(x,y)=x+2y$, calculate the slope of an indifference curve and draw two curves.
\end{enumerate}

\subsection*{Analytical exercises}
\begin{enumerate}[resume,label=\arabic*.]
  \item Let $u(x,y)=x^{1/3}y^{2/3}$. Calculate $MU_x$, $MU_y$, and $MRS_{x,y}$.
  \item A consumer has utility $u(x,y)=xy+2x$. Find the equation of the indifference curve through $(3,4)$.
  \item For $u(x,y)=\min\{2x,y\}$, describe the bundles on the indifference curve associated with utility $10$.
\end{enumerate}

\subsection*{Exam-style multiple-choice questions}
\begin{enumerate}[resume,label=\arabic*.]
  \item If two utility functions rank every pair of bundles in the same order, then: (a) they represent different preferences; (b) they represent the same ordinal preferences; (c) they must assign identical utility values; or (d) neither can be used in consumer theory.
  \item If $u(x,y)=xy$, the marginal utility of $x$ is: (a) $x$; (b) $y$; (c) $xy$; or (d) $y/x$.
  \item Perfect-complement preferences are represented by indifference curves that are: (a) straight lines; (b) upward sloping; (c) L-shaped; or (d) horizontal.
\end{enumerate}

\subsection*{PPE discussion question}

\noindent Does an ordinal account of preferences provide enough information to compare social outcomes? Discuss the distinction between identifying a person's own ranking of options and making interpersonal comparisons of well-being.
